\documentclass[11pt]{article}
\usepackage[margin=1in]{geometry}
\usepackage{amsmath,amssymb}
\usepackage{parskip}
\usepackage{enumitem}
\usepackage[colorlinks=true,linkcolor=blue,urlcolor=blue]{hyperref}

\newcommand{\term}[1]{\medskip\noindent\textbf{#1.}\ }

\title{\vspace{-1.2cm}A plain-language glossary for crystal symmetry\\[2pt]
in ISDF--THC (diamond as the running example)}
\author{}
\date{\today}

\begin{document}
\maketitle
\vspace{-0.8cm}

This is a short dictionary of the symmetry words we have been using --- \emph{star}, \emph{little
group}, \emph{high symmetry}, \emph{orbit}, \emph{copy}, and so on. Each entry is one plain sentence,
a quick picture, and a diamond number where it helps. Read it top to bottom: the terms build on each
other. Diamond has $|P|=48$ point-group operations and $n_{\mathrm{IP}}=136$ interpolation points per
cell.

\section*{1.\ The crystal and its symmetries}

\term{Unit cell / lattice} The unit cell is the smallest repeating box of atoms; copying it on a
regular grid (the lattice) builds the whole crystal.

\term{Space group} The complete list of moves that leave the crystal looking identical --- rotations,
reflections, and translations, and combinations of them. Diamond's is called $Fd\bar{3}m$.

\term{Point group} The moves that hold one point fixed (rotations and reflections, no sliding). Diamond
has $|P|=48$ of them. Think of it as the symmetry of a single site.

\term{Symmorphic vs.\ nonsymmorphic} Symmorphic means every symmetry is a clean rotation/reflection
(plus a whole-cell shift). Nonsymmorphic means some symmetries need a \emph{fractional} shift --- a
``glide'' or ``screw.'' Diamond is nonsymmorphic. In practice this only sprinkles in extra phases;
none of the ideas below change.

\section*{2.\ Momentum space}

\term{$k$-point (momentum)} A label for how a wave's phase steps from one cell to the next. Rather than
follow every cell, we follow a momentum $k$. There are $N_k$ of them.

\term{BvK supercell} Born--von K\'arm\'an: a finite box of $N_k$ cells with wrap-around (periodic)
edges. $N_k$ = number of cells = number of $k$-points (8 for a $2\times2\times2$ mesh, 216 for
$6\times6\times6$).

\term{Brillouin zone} The space of allowed $k$'s. A few special spots have names --- $\Gamma$ (the
center), $X$, $L$.

\section*{3.\ Grouping the $k$-points (the heart of it)}

\term{Star (a.k.a.\ family)} Take one $k$ and hit it with all 48 point-group moves. You get a
\emph{set} of $k$'s --- that is its star. Every $k$ in a star is a symmetry-copy of the others, so you
only need to study \emph{one} of them. A star has at most 48 members.

\term{Little group (of $k$)} The subset of the 48 moves that leave \emph{this particular} $k$ where it
is (up to a whole reciprocal-lattice jump). Rule of thumb (orbit--stabilizer):
\[
   (\text{star size})\times(\text{little-group size}) = 48 .
\]
A large little group means few star members; a small little group means many.

\term{High-symmetry vs.\ generic (low-symmetry) $k$} A \emph{high-symmetry} $k$ has a \emph{big} little
group (lots of moves fix it) and therefore a \emph{small} star. A \emph{generic / low-symmetry} $k$ has
a \emph{small} little group and a \emph{big} star (up to 48). $\Gamma$ is the most symmetric point ---
its little group is all 48 moves. Coarse $k$-meshes are made almost entirely of high-symmetry points;
as you refine the mesh, the new points are mostly low-symmetry.

\section*{4.\ The interpolation grid}

\term{Interpolation points (ISDF), $n_{\mathrm{IP}}$} ISDF represents products of orbitals by their
values at a small, cleverly chosen set of real-space points --- the interpolation points. Diamond uses
$n_{\mathrm{IP}}=136$ per cell.

\term{Fiber} The bundle of $n_{\mathrm{IP}}$ interpolation points sitting ``over'' one $k$. The whole
grid is $N_k$ fibers $\times\ n_{\mathrm{IP}}$ points $=N_k\cdot n_{\mathrm{IP}}$ in total.

\term{Orbit (of interpolation points)} Fix one $k$ and apply its little-group moves to one
interpolation point: you sweep out a \emph{set} of points --- an orbit. Points only ever mix
\emph{within} an orbit, never across orbits. An orbit is at most as big as the little group. (Do not
confuse this with a star, which is an orbit of \emph{$k$-points}; here it is an orbit of
\emph{interpolation points}.)

\section*{5.\ Representation-theory words}

\term{Representation} A rule that turns each symmetry move into a matrix acting on some space --- here,
how the moves shuffle-and-phase the interpolation points.

\term{Irreducible representation (irrep)} A representation that cannot be split into smaller
independent pieces --- the ``atoms'' of symmetry. Each irrep has a fixed size, its \emph{dimension}.

\term{Dimension and partner ($d_\mu$, $\nu$)} An irrep of dimension $d_\mu$ comes with $d_\mu$
\emph{partners} ($\nu=1,\dots,d_\mu$) that the symmetry moves rotate into one another --- like the
$(x,y)$ pair under rotation. Diamond's little-group irreps have dimensions 1, 2, or 3.

\term{Multiplicity and copy ($n_\mu$, $c$)} How many separate times a given irrep appears in a space;
each appearance is a \emph{copy}, labelled $c$. Symmetry cannot tell the copies apart --- so this is
exactly where the real, material-specific numbers live.

\term{Character} The trace (sum of the diagonal) of an irrep's matrices. A basis-independent
\emph{fingerprint} that identifies the irrep.

\term{Character projection} A formula built from characters that automatically extracts the part of a
space belonging to a chosen irrep. It is the non-commutative cousin of a Fourier filter.

\term{Schur's lemma / commutant} The theorem behind all the savings: any operator that respects the
symmetry must be \emph{identity on the partners} and \emph{free on the copies}, and block-diagonal by
irrep. That is precisely why symmetry compresses the data.

\section*{6.\ Putting it together: the transform and the blocks}

\term{$Q$ (symmetry adaptation)} The change of basis from the raw labels
$|k\rangle|I\rangle$ (momentum, interpolation point) to the symmetry labels
$|i\rangle|a\rangle|c\rangle$. It is the ``generalized Fourier transform'' of the whole space group
(ordinary Fourier on the translations, character projection on the point group).

\term{The symmetry labels $(i,a,c)$}
\begin{itemize}[leftmargin=1.6em,topsep=2pt]
  \item $i=(f,\mu)$: which \emph{star} $f$ and which little-group \emph{irrep} $\mu$ --- together, the
        full-crystal irrep.
  \item $a=(\kappa,\nu)$: which \emph{member} $\kappa$ of the star, and which \emph{partner} $\nu$ of
        the irrep --- together, the partner index (size $m_i=|\text{star}|\cdot d_\mu$).
  \item $c$: the \emph{copy} --- the multiplicity slot, where the material data sits.
\end{itemize}

\term{$W$ block structure} In the symmetry basis the Coulomb tensor is
$W=\bigoplus_i \mathbb{1}_{m_i}\otimes w_i$: \emph{block-diagonal by} $i=(f,\mu)$, \emph{identity on the
partner} $a=(\kappa,\nu)$ (so all star members and all partners share one block), and a \emph{free
matrix} $w_i$ only on the copies $c$.

\term{Fold-back lattice vector ($L$)} When a symmetry move carries an interpolation point out of its
home cell into a neighbour, $L$ records which neighbour. It supplies the Bloch phase $e^{-ik\cdot L}$
that relates one star member to another.

\term{Unique data} The count of genuinely independent numbers left after symmetry:
$\sum_{\text{orbits}}|o|^2$ for the transform ($6800$ at $2\times2\times2$), $\sum_i n_i^2$ for $W$
($3384$). The remarkable part: these stay almost \emph{flat} as $N_k$ grows (the added low-symmetry
points contribute tiny blocks), even though $N_k$ itself grows fast.

\end{document}
